\documentclass[11pt]{article}

\usepackage[margin=1in]{geometry}
\usepackage{amsmath,amssymb,amsthm,mathtools}
\usepackage{enumitem}
\usepackage{hyperref}
\usepackage{xcolor}

\hypersetup{
  colorlinks=true,
  linkcolor=blue!60!black,
  citecolor=blue!60!black,
  urlcolor=blue!60!black
}

\newtheorem{theorem}{Theorem}
\newtheorem{lemma}{Lemma}
\newtheorem{proposition}{Proposition}
\newtheorem{corollary}{Corollary}
\newtheorem{assumption}{Assumption}
\newtheorem{remark}{Remark}
\newcommand{\E}{\mathbb{E}}
\newcommand{\Prb}{\mathbb{P}}
\newcommand{\ind}{\mathbf{1}}
\newcommand{\calA}{\mathcal{A}}
\newcommand{\calF}{\mathcal{F}}

\title{Selfish Smooth Bidding for Target Scheduling:\\A Completion-Rate Bound with Private Urgency}
\author{}
\date{}

\begin{document}
\maketitle

\section{Model}

We consider a finite-horizon target-scheduling problem. At each decision epoch the agent observes a state
$s$ and an active target set
\begin{equation}
  A(s) \subseteq \{1,\dots,k\}, \qquad |A(s)| = n(s) \le k .
  \label{eq:active-set}
\end{equation}
At most one target is selected at a decision epoch. If target $i$ is selected, the agent follows the local policy
for target $i$ until that attempt terminates, either by successfully reaching the target or by consuming time and
returning to the next decision state. New targets may appear after old targets are reached, but the active set size
always satisfies \eqref{eq:active-set}.

Let $T$ denote the horizon. Let
\begin{equation}
  N^{\mathrm{NE}}(T)
\end{equation}
be the number of successful target completions under the bidding-induced Nash scheduler, and let
\begin{equation}
  N^\star(T)
\end{equation}
be the maximum number of completions achievable by an optimal scheduler.

\section{Smooth bidding mechanism}

At state $s$, each active target-player $i\in A(s)$ submits a bid
\begin{equation}
  b_i \in [0,B] .
  \label{eq:bid-domain}
\end{equation}
Given a bid vector $b$, target $i$ is selected with probability
\begin{equation}
  p_i(b;s)
  =
  \frac{\eta+b_i}{n(s)\eta + \sum_{j\in A(s)} b_j},
  \qquad \eta>0.
  \label{eq:proportional-allocation}
\end{equation}
The offset $\eta$ ensures that if all bids are zero, the scheduler chooses uniformly among active targets.

\subsection{Cross-reachability and private urgency}

The key distinction is between raw reachability value and urgency. Define
\begin{equation}
  r_{ij}(s)
  =
  \Prb\{\text{target }i\text{ is eventually reached if target }j\text{ is selected now}\mid s\}.
  \label{eq:cross-reachability}
\end{equation}
Thus
\begin{equation}
  q_i(s) := r_{ii}(s)
  \label{eq:self-success}
\end{equation}
is the success probability for target $i$ if it is selected now. The off-diagonal term $r_{ij}(s)$ measures
how much selecting target $j$ helps target $i$.

Given allocation probabilities $p=(p_j)_{j\in A(s)}$, the expected reachability payoff of player $i$ is
\begin{equation}
  R_i(s,p)
  =
  \sum_{j\in A(s)} p_j r_{ij}(s).
  \label{eq:expected-private-reachability}
\end{equation}
The private urgency of player $i$ is the marginal gain from increasing the probability of selecting $i$ rather than
using the current randomized allocation:
\begin{equation}
  u_i(s,p)
  =
  q_i(s)-R_i(s,p)
  =
  r_{ii}(s)-\sum_{j\in A(s)}p_jr_{ij}(s).
  \label{eq:urgency}
\end{equation}
If nearby targets help player $i$ remain reachable, then $R_i(s,p)$ is large and $u_i(s,p)$ is small. Thus
\eqref{eq:urgency} captures selfish free-riding on nearby objectives without introducing any global coordination reward.

For the smooth bidding game, player $i$ has one-step utility
\begin{equation}
  U_i(b;s)
  =
  \sum_{j\in A(s)} p_j(b;s) r_{ij}(s) - \rho b_i.
  \label{eq:player-utility}
\end{equation}
Let
\begin{equation}
  x_i=\eta+b_i,
  \qquad
  S=\sum_{j\in A(s)}x_j.
  \label{eq:x-and-S}
\end{equation}
Then $p_i=x_i/S$. A direct derivative calculation gives
\begin{equation}
  \frac{\partial U_i}{\partial b_i}(b;s)
  =
  \frac{1}{S}\left(r_{ii}(s)-\sum_{j\in A(s)}p_j(b;s)r_{ij}(s)\right)-\rho
  =
  \frac{u_i(s,p(b;s))}{S}-\rho .
  \label{eq:marginal-utility}
\end{equation}
Thus a player bids until private urgency equals the marginal bid cost, subject to the bid constraints in
\eqref{eq:bid-domain}. This is the formal justification for treating bids as urgency-sensitive rather than merely
value-sensitive.

\section{Mechanism quality}

Let $p^{\mathrm{NE}}(s)$ be the selection distribution induced by a Nash equilibrium of the one-step bidding game
at state $s$. The following result is purely algebraic and does not require any global coordination rewards.

\subsection{Urgency mass and the parameter $\gamma$}

Define
\begin{equation}
  u_{\max}(s)
  =
  \max_{i\in A(s)} u_i(s,p^{\mathrm{NE}}(s)).
  \label{eq:u-max}
\end{equation}
For parameters $\lambda\in(0,1]$ and a state $s$, define the high-urgency set
\begin{equation}
  H_\lambda(s)
  =
  \{i\in A(s): u_i(s,p^{\mathrm{NE}}(s))\ge \lambda u_{\max}(s)\}.
  \label{eq:high-urgency-set}
\end{equation}
Let
\begin{equation}
  \beta(s)
  =
  \sum_{i\in H_\lambda(s)} p_i^{\mathrm{NE}}(s)
  \label{eq:beta-def}
\end{equation}
be the probability mass that the equilibrium scheduler places on high-urgency targets.

Define the urgency quality of the allocation by
\begin{equation}
  \gamma_u(s)
  =
  \frac{\sum_{i\in A(s)} p_i^{\mathrm{NE}}(s)u_i(s,p^{\mathrm{NE}}(s))}{u_{\max}(s)} .
  \label{eq:gamma-u}
\end{equation}

\begin{lemma}[Urgency-mass lower bound]
If $\beta(s)$ is defined by \eqref{eq:beta-def}, then
\begin{equation}
  \gamma_u(s) \ge \lambda \beta(s).
  \label{eq:gamma-urgency-bound}
\end{equation}
\end{lemma}

\begin{proof}
By definition of $H_\lambda(s)$ in \eqref{eq:high-urgency-set}, every $i\in H_\lambda(s)$ satisfies
$u_i(s,p^{\mathrm{NE}}(s))\ge \lambda u_{\max}(s)$. Therefore
\begin{align}
  \sum_{i\in A(s)} p_i^{\mathrm{NE}}(s)u_i(s,p^{\mathrm{NE}}(s))
  &\ge
  \sum_{i\in H_\lambda(s)} p_i^{\mathrm{NE}}(s)u_i(s,p^{\mathrm{NE}}(s)) \notag\\
  &\ge
  \lambda u_{\max}(s)
  \sum_{i\in H_\lambda(s)}p_i^{\mathrm{NE}}(s) \notag\\
  &=
  \lambda u_{\max}(s)\beta(s). \label{eq:gamma-proof-step}
\end{align}
Dividing \eqref{eq:gamma-proof-step} by $u_{\max}(s)$ gives \eqref{eq:gamma-urgency-bound}.
\end{proof}

\subsection{Bridge from urgency to success probability}

The completion count depends on actual success probabilities $q_i(s)$, not on urgency alone. Since
\begin{equation}
  u_i(s,p) = q_i(s)-R_i(s,p) \le q_i(s),
  \label{eq:urgency-below-success}
\end{equation}
urgency is a conservative proxy for success. To compare urgency quality with success quality, we use the following
bridge parameter.

\begin{assumption}[Urgency-success bridge]
There exists $\kappa\in[0,1]$ such that for every reachable state $s$,
\begin{equation}
  u_{\max}(s) \ge \kappa q_{\max}(s),
  \qquad
  q_{\max}(s)=\max_{i\in A(s)}q_i(s).
  \label{eq:bridge-assumption}
\end{equation}
\end{assumption}

The bridge parameter $\kappa$ is positive when the most valuable target is also genuinely urgent: selecting other
targets does not already give almost the same chance of eventually reaching it. When there are no deadlines and no
horizon pressure, $\kappa$ may be close to zero, and urgency-based prioritization is not informative.

\begin{lemma}[From urgency mass to success mass]
Suppose \eqref{eq:bridge-assumption} holds and $\beta(s)\ge \beta$ uniformly over reachable states. Then
\begin{equation}
  \sum_{i\in A(s)}p_i^{\mathrm{NE}}(s)q_i(s)
  \ge
  \lambda\beta\kappa\, q_{\max}(s).
  \label{eq:success-gamma-bound}
\end{equation}
Equivalently, the success-probability quality
\begin{equation}
  \gamma_q(s)
  =
  \frac{\sum_i p_i^{\mathrm{NE}}(s)q_i(s)}{q_{\max}(s)}
  \label{eq:gamma-q-def}
\end{equation}
satisfies
\begin{equation}
  \gamma_q(s)\ge \lambda\beta\kappa.
  \label{eq:gamma-q-lower}
\end{equation}
\end{lemma}

\begin{proof}
By \eqref{eq:urgency-below-success}, $q_i(s)\ge u_i(s,p^{\mathrm{NE}}(s))$ for each $i$. Therefore
\begin{align}
  \sum_i p_i^{\mathrm{NE}}(s)q_i(s)
  &\ge
  \sum_i p_i^{\mathrm{NE}}(s)u_i(s,p^{\mathrm{NE}}(s)) \notag\\
  &\ge
  \lambda\beta(s)u_{\max}(s) \notag\\
  &\ge
  \lambda\beta\kappa q_{\max}(s), \label{eq:success-gamma-proof}
\end{align}
where the second inequality uses Lemma 1 and the third uses \eqref{eq:bridge-assumption} and $\beta(s)\ge\beta$.
This proves \eqref{eq:success-gamma-bound} and \eqref{eq:gamma-q-lower}.
\end{proof}

\section{Completion-rate approximation theorem}

The factor $\underline\mu/\overline\mu$ is a rate comparison. The Nash scheduler's completion rate is controlled by
its success probability per attempt divided by its expected attempt duration. The optimal scheduler is upper bounded
by the fastest possible expected time per successful completion.

Let $\tau_i(s)$ be the random time consumed by selecting target $i$ at state $s$. Let $Y$ be the indicator that the
attempt succeeds.

\begin{assumption}[Expected time regularity]
There are constants $0<\underline\mu\le \overline\mu<\infty$ such that:
\begin{enumerate}[label=(\roman*)]
  \item for every reachable state $s$ and every target selected by the Nash scheduler,
  \begin{equation}
    \E[\tau_i(s)\mid s] \le \overline\mu;
    \label{eq:upper-mu}
  \end{equation}
  \item no scheduler can obtain successful completions at expected time per success smaller than $\underline\mu$, i.e.
  \begin{equation}
    \E[N^\star(T)]\le \frac{T}{\underline\mu}.
    \label{eq:optimal-rate-upper}
  \end{equation}
\end{enumerate}
\end{assumption}

The second condition can be derived, for example, when every successful completion has conditional expected duration
at least $\underline\mu$ under any policy, with the usual optional-stopping regularity conditions.

\begin{assumption}[Attempt-rate regularity]
There is a finite overshoot constant $C_{\mathrm{os}}$ such that if attempts under the Nash scheduler have conditional
expected duration at most $\overline\mu$, then the expected number of attempts started before time $T$ is at least
\begin{equation}
  \E[M^{\mathrm{NE}}(T)] \ge \frac{T}{\overline\mu}-C_{\mathrm{os}} .
  \label{eq:attempt-count-lower}
\end{equation}
This holds, for instance, under standard renewal-type conditions with bounded second moments; if attempt durations are
uniformly bounded, $C_{\mathrm{os}}$ can be taken as a constant independent of $T$.
\end{assumption}

\begin{assumption}[Availability]
There exists $q_0\in(0,1]$ such that at every nonterminal reachable state,
\begin{equation}
  q_{\max}(s)\ge q_0.
  \label{eq:q0-assumption}
\end{equation}
\end{assumption}

\begin{theorem}[Expected completion-rate bound]\label{thm:completion-rate-bound}
Suppose Assumptions 1--3 hold. Suppose also that the equilibrium allocation satisfies
\begin{equation}
  \beta(s)\ge \beta
  \label{eq:beta-uniform}
\end{equation}
for the high-urgency set in \eqref{eq:high-urgency-set}. Then
\begin{equation}
  \E[N^{\mathrm{NE}}(T)]
  \ge
  \lambda\beta\kappa q_0
  \left(\frac{T}{\overline\mu}-C_{\mathrm{os}}\right).
  \label{eq:ne-completions-lower}
\end{equation}
Moreover,
\begin{equation}
  \frac{\E[N^{\mathrm{NE}}(T)]}{\E[N^\star(T)]}
  \ge
  \lambda\beta\kappa q_0
  \left(
    \frac{\underline\mu}{\overline\mu}
    -
    \frac{C_{\mathrm{os}}\underline\mu}{T}
  \right).
  \label{eq:main-ratio-bound}
\end{equation}
In particular, for large horizons,
\begin{equation}
  \E[N^{\mathrm{NE}}(T)]
  \gtrsim
  \lambda\beta\kappa q_0
  \frac{\underline\mu}{\overline\mu}
  \E[N^\star(T)].
  \label{eq:asymptotic-bound}
\end{equation}
\end{theorem}

\begin{proof}
At any decision state $s$, Lemma 2 and \eqref{eq:q0-assumption} imply that the conditional success probability of
the Nash-selected attempt is at least
\begin{equation}
  \sum_i p_i^{\mathrm{NE}}(s)q_i(s)
  \ge
  \lambda\beta\kappa q_{\max}(s)
  \ge
  \lambda\beta\kappa q_0 .
  \label{eq:conditional-success-lower}
\end{equation}
Let $M^{\mathrm{NE}}(T)$ be the number of attempts started before time $T$. Summing the conditional success lower bound
\eqref{eq:conditional-success-lower} over attempts and using \eqref{eq:attempt-count-lower} gives
\begin{equation}
  \E[N^{\mathrm{NE}}(T)]
  \ge
  \lambda\beta\kappa q_0\, \E[M^{\mathrm{NE}}(T)]
  \ge
  \lambda\beta\kappa q_0
  \left(\frac{T}{\overline\mu}-C_{\mathrm{os}}\right),
  \label{eq:ne-bound-proof}
\end{equation}
which proves \eqref{eq:ne-completions-lower}.

By \eqref{eq:optimal-rate-upper},
\begin{equation}
  \E[N^\star(T)]\le \frac{T}{\underline\mu}.
  \label{eq:opt-bound-proof}
\end{equation}
Combining \eqref{eq:ne-bound-proof} and \eqref{eq:opt-bound-proof} yields
\begin{align}
  \frac{\E[N^{\mathrm{NE}}(T)]}{\E[N^\star(T)]}
  &\ge
  \lambda\beta\kappa q_0
  \left(\frac{T}{\overline\mu}-C_{\mathrm{os}}\right)
  \frac{\underline\mu}{T} \notag\\
  &=
  \lambda\beta\kappa q_0
  \left(
    \frac{\underline\mu}{\overline\mu}
    -
    \frac{C_{\mathrm{os}}\underline\mu}{T}
  \right),
  \label{eq:ratio-proof-step}
\end{align}
which is \eqref{eq:main-ratio-bound}. The asymptotic form \eqref{eq:asymptotic-bound} follows by suppressing the
$O(1/T)$ overshoot term.
\end{proof}

\begin{remark}[Special case: deterministic feasible completions]
If every feasible selected target is reached with probability one, then $q_i(s)=1$ for every feasible active target.
In that case $q_0=1$. If there is meaningful deadline or horizon pressure, then $\kappa$ is positive. If there is no
such pressure, urgencies may be near zero even though all $q_i=1$, and the urgency-based part of the theorem becomes
vacuous; the scheduling problem is then governed mostly by route-efficiency, not deadline urgency.
\end{remark}

\section{Estimating the urgency fraction $\alpha$}

The theorem is stated in terms of $\beta$, the probability mass placed on high-urgency targets. It is often useful to
first estimate the fraction of high-urgency targets,
\begin{equation}
  \alpha(s)=\frac{|H_\lambda(s)|}{|A(s)|}.
  \label{eq:alpha-state}
\end{equation}
The relationship between $\alpha$ and $\beta$ depends on how strongly bidding concentrates probability on urgent
targets. If high-urgency targets bid at least $b_H$ and the remaining targets bid at most $b_L$, then by
\eqref{eq:proportional-allocation}, with $h=|H_\lambda(s)|/|A(s)|$,
\begin{equation}
  \beta(s)
  \ge
  \frac{h(\eta+b_H)}{h(\eta+b_H)+(1-h)(\eta+b_L)}.
  \label{eq:beta-from-alpha}
\end{equation}
Thus sparse urgent targets can still receive large probability mass if their bids separate from low-urgency targets.
This is where bidding can improve over uniform random selection, for which $\beta_{\mathrm{rand}}(s)=h$.

\subsection{Fixed target lifetimes}

Suppose each target remains alive for a fixed lifetime $L$ after it appears. If target $i$ appeared at time $a_i^{\rm app}$,
then its remaining lifetime, or slack, at time $t$ is
\begin{equation}
  h_i(t)=L-(t-a_i^{\rm app}).
  \label{eq:slack-def}
\end{equation}
Let $\mu$ be a typical completion-time scale. For a threshold multiplier $c>0$, define a deadline-urgent target by
\begin{equation}
  h_i(t)\le c\mu.
  \label{eq:deadline-urgent}
\end{equation}
Assume active target ages are approximately independent and uniformly distributed on $[0,L]$, as in a stationary
asynchronous arrival or respawn process. Then
\begin{equation}
  \Prb\{h_i(t)\le c\mu\}
  =
  \min\left\{1,\frac{c\mu}{L}\right\}.
  \label{eq:alpha-fixed-lifetime}
\end{equation}
Therefore the expected fraction of deadline-urgent targets is
\begin{equation}
  \alpha
  =
  \min\left\{1,\frac{c\mu}{L}\right\}.
  \label{eq:alpha-final}
\end{equation}
If $n=|A(s)|$ active target ages are independent, then
\begin{equation}
  |H|\sim \mathrm{Binomial}(n,\alpha).
  \label{eq:binomial-H}
\end{equation}
By Hoeffding's inequality, with probability at least $1-\delta$,
\begin{equation}
  \left|\frac{|H|}{n}-\alpha\right|
  \le
  \sqrt{\frac{\log(2/\delta)}{2n}}.
  \label{eq:hoeffding-alpha}
\end{equation}
Thus $\alpha$ is estimable from $L$ and $\mu$, with concentration improving as the number of active targets grows.

\subsection{No-expiry targets}

If targets never expire, then $L=\infty$ and their deadline slack is infinite. Under the deadline-urgency definition
\eqref{eq:deadline-urgent},
\begin{equation}
  \alpha = 0.
  \label{eq:no-expiry-alpha}
\end{equation}
The urgency theory is then vacuous unless urgency is redefined relative to the global horizon. For example, one can
replace $h_i(t)$ with the remaining episode time
\begin{equation}
  R_t=T-t
  \label{eq:remaining-horizon}
\end{equation}
and call a target horizon-urgent if delaying it makes completion before $T$ unlikely. Without expiries or horizon
pressure, all unreached targets with $q_i=1$ are equally valuable to their owners, and selfish bidding has little
reason to produce nontrivial prioritization.


\section{Hard winner-pays bidding}
\label{sec:winner-pays}

The smooth mechanism in \eqref{eq:proportional-allocation} gives a differentiable way to translate urgency into
selection probabilities. A hard winner-take-all auction gives a different type of result: it can \emph{implement} the
maximum-urgency target as a Nash winner when the bid grid and winner payment are calibrated to the urgency gap.
This is an equilibrium-selection result, not a universal guarantee over all possible equilibria.

\subsection{One-shot urgency contest}
Fix a decision state $s$ and write
\begin{equation}
  u_i = u_i(s,p)
  \label{eq:hard-u-i}
\end{equation}
for the private marginal urgency of target $i$. In this subsection the quantities $u_i$ are treated as fixed urgency
scores for the current decision. Let
\begin{equation}
  u_{(1)}\ge u_{(2)}\ge \cdots \ge u_{(n)}
  \label{eq:sorted-urgencies}
\end{equation}
be the sorted urgencies, and let $i^\star$ be a target with urgency $u_{(1)}$.

Bids lie on a discrete grid
\begin{equation}
  \mathcal B=\{0,\Delta,2\Delta,\dots,B\}.
  \label{eq:hard-bid-grid}
\end{equation}
The target with the largest bid wins; ties are broken uniformly at random among the highest bidders. Under a
winner-pays rule, only the winner pays. Normalizing away the baseline reachability payoff that a player receives when
it does not win, the incremental payoff of player $i$ is
\begin{equation}
  U_i^{\mathrm{WP}}(b)
  =
  \Prb\{i\text{ wins under }b\}\,\bigl(u_i-\rho b_i\bigr).
  \label{eq:winner-pays-utility}
\end{equation}
Thus $u_i$ is the private value of winning control at the current decision, and $\rho b_i$ is the payment made only if
player $i$ wins.

\begin{proposition}[Winner-pays implementation of the maximum-urgency target]
\label{prop:winner-pays-implementation}
Suppose there exists a grid bid $b^\star\in\mathcal B$ with $b^\star\ge \Delta$ such that
\begin{equation}
  u_{(2)} \le \rho b^\star
  \label{eq:wp-deter-losers}
\end{equation}
and
\begin{equation}
  \rho(b^\star+\Delta) \le u_{(1)}.
  \label{eq:wp-no-lower}
\end{equation}
Then there exists a pure Nash equilibrium of the winner-pays hard auction in which $i^\star$ wins with probability
one. In particular, the hard auction implements the maximum-urgency decision.
\end{proposition}

\begin{proof}
Consider the following bid profile. Player $i^\star$ bids $b^\star$. One runner-up with urgency at most $u_{(2)}$ bids
$b^\star-\Delta$, and all remaining players bid at most $b^\star-\Delta$; if $b^\star=\Delta$, this means the
runner-up bids zero. Thus $i^\star$ wins uniquely and obtains payoff
\begin{equation}
  U_{i^\star}^{\mathrm{WP}}
  =
  u_{(1)}-\rho b^\star.
  \label{eq:wp-top-payoff}
\end{equation}

First consider a losing player $j\ne i^\star$. Any bid below $b^\star$ still loses and gives payoff zero. A deviation
to $b^\star$ creates a tie and yields payoff at most
\begin{equation}
  \frac{1}{2}\bigl(u_j-\rho b^\star\bigr),
  \label{eq:wp-loser-tie-payoff}
\end{equation}
which is nonpositive by \eqref{eq:wp-deter-losers}. Any deviation above $b^\star$ costs at least
$\rho(b^\star+\Delta)\ge \rho b^\star\ge u_{(2)}\ge u_j$ and is also nonprofitable. Hence no loser has a profitable
deviation.

Now consider player $i^\star$. Raising its bid only increases the payment while preserving the same winning
probability, so it is not profitable. Lowering below $b^\star-\Delta$ loses or weakly lowers the probability of
winning. The only potentially profitable downward deviation is to $b^\star-\Delta$, which creates a tie with the
runner-up and gives payoff at most
\begin{equation}
  \frac{1}{2}\bigl(u_{(1)}-\rho(b^\star-\Delta)\bigr).
  \label{eq:wp-top-lower-payoff}
\end{equation}
The original winning payoff in \eqref{eq:wp-top-payoff} is at least \eqref{eq:wp-top-lower-payoff} whenever
\begin{equation}
  u_{(1)}-\rho b^\star
  \ge
  \frac{1}{2}\bigl(u_{(1)}-\rho(b^\star-\Delta)\bigr),
  \label{eq:wp-no-lower-derive}
\end{equation}
which is equivalent to
\begin{equation}
  \rho(b^\star+\Delta)\le u_{(1)}.
  \label{eq:wp-no-lower-equivalent}
\end{equation}
This is exactly \eqref{eq:wp-no-lower}. Therefore $i^\star$ also has no profitable deviation. The profile is a pure
Nash equilibrium and $i^\star$ wins uniquely.
\end{proof}

\begin{remark}[Interpretation of the bid-separation condition]
The two conditions \eqref{eq:wp-deter-losers}--\eqref{eq:wp-no-lower} require a bid level whose cost lies between the
second-largest urgency and the largest urgency, with one grid step of slack:
\begin{equation}
  u_{(2)} \le \rho b^\star \le \rho(b^\star+\Delta) \le u_{(1)}.
  \label{eq:wp-gap-interpretation}
\end{equation}
Thus the result is strongest when the top urgency is separated from the rest and the bid grid is fine enough that
$\rho\Delta$ is small compared with the urgency gap. If the top two urgencies are nearly tied, the hard auction may
have multiple equilibria or tie-based equilibria; the implementation result then applies to an equilibrium selection,
not to every equilibrium.
\end{remark}

\subsection{Completion guarantee for the implemented winner-pays equilibrium}

Proposition~\ref{prop:winner-pays-implementation} implies that, at any state where the bid-separation condition holds,
the hard winner-pays auction can select a maximum-urgency target. In the notation of Section~3, this means that the
high-urgency set with $\lambda=1$ receives probability mass
\begin{equation}
  \beta(s)=1.
  \label{eq:wp-beta-one}
\end{equation}
Combining \eqref{eq:wp-beta-one} with Theorem~\ref{thm:completion-rate-bound} gives the following corollary.

\begin{corollary}[Winner-pays completion-rate bound]
\label{cor:winner-pays-completion}
Suppose that at every reachable decision state the winner-pays hard auction satisfies the bid-separation conditions
\eqref{eq:wp-deter-losers}--\eqref{eq:wp-no-lower} and the equilibrium from Proposition~\ref{prop:winner-pays-implementation}
is selected. Suppose also that the urgency-success bridge, expected time regularity, attempt-rate regularity, and
availability assumptions hold. Then
\begin{equation}
  \E[N^{\mathrm{WP}}(T)]
  \ge
  \kappa q_0
  \left(\frac{T}{\overline\mu}-C_{\mathrm{os}}\right),
  \label{eq:wp-ne-completions-lower}
\end{equation}
and
\begin{equation}
  \frac{\E[N^{\mathrm{WP}}(T)]}{\E[N^\star(T)]}
  \ge
  \kappa q_0
  \left(
    \frac{\underline\mu}{\overline\mu}
    -
    \frac{C_{\mathrm{os}}\underline\mu}{T}
  \right).
  \label{eq:wp-main-ratio-bound}
\end{equation}
For large horizons,
\begin{equation}
  \E[N^{\mathrm{WP}}(T)]
  \gtrsim
  \kappa q_0\frac{\underline\mu}{\overline\mu}\E[N^\star(T)].
  \label{eq:wp-asymptotic-bound}
\end{equation}
\end{corollary}

\begin{proof}
Under the selected equilibrium from Proposition~\ref{prop:winner-pays-implementation}, the winner has maximum urgency
at each decision state. Therefore in Theorem~\ref{thm:completion-rate-bound} we can take $\lambda=1$ and $\beta=1$.
Substituting these values into \eqref{eq:ne-completions-lower} and \eqref{eq:main-ratio-bound} gives
\eqref{eq:wp-ne-completions-lower} and \eqref{eq:wp-main-ratio-bound}. The asymptotic form follows as in
\eqref{eq:asymptotic-bound}.
\end{proof}

\begin{remark}[Why this is not a universal hard-auction guarantee]
The winner-pays result is an implementation theorem. It says that, when a grid bid can separate the highest urgency
from the rest, there exists a Nash equilibrium that implements the maximum-urgency target. It does not imply that all
winner-pays equilibria are optimal, nor does it remove the need for the completion-time and urgency-success bridge
conditions. For all-pay hard auctions the corresponding statement is typically weaker, because losing bidders also
pay; one usually obtains threshold or mixed-equilibrium guarantees rather than the clean pure implementation result
above.
\end{remark}

\section{Interpretation}

The bound in \eqref{eq:main-ratio-bound} decomposes the performance of the selfish smooth bidding scheduler into four
terms:
\begin{equation}
  \lambda\beta\kappa q_0\cdot \frac{\underline\mu}{\overline\mu}.
  \label{eq:factor-decomposition}
\end{equation}
The terms have the following meanings:
\begin{itemize}[leftmargin=2em]
  \item $\lambda$: how urgent the high-urgency set is relative to the maximum urgency;
  \item $\beta$: how much probability mass the Nash bidding mechanism places on that high-urgency set;
  \item $\kappa$: how much maximum urgency reflects maximum actual success probability;
  \item $q_0$: a lower bound on the best available success probability;
  \item $\underline\mu/\overline\mu$: the completion-rate regularity factor, comparing the fastest possible expected time per success with the Nash scheduler's expected time per attempt.
\end{itemize}
This formulation avoids identifying urgency with value. Raw success is $q_i=r_{ii}$, while urgency is the marginal
quantity $u_i=q_i-\sum_jp_jr_{ij}$. Bidding is useful precisely when the equilibrium concentrates probability mass on
objectives with large private marginal urgency.

\end{document}
